\documentclass[11pt]{article}
\usepackage[a4paper,margin=1in]{geometry}
\usepackage{amsmath,amssymb,amsthm}
\usepackage{booktabs}
\usepackage{array}
\usepackage{hyperref}

\newtheorem{theorem}{Theorem}
\newtheorem{conjecture}{Conjecture}

\newcommand{\CRE}{\mathrm{CRE}}
\newcommand{\CRCFT}{\mathrm{CRCFT}}
\newcommand{\CTMT}{\mathrm{CTMT}}
\newcommand{\XiC}{\Xi_C}

\title{Step 187: Framework RH-Carrier Dichotomy Theorem}
\author{RH Membrane Adequacy Track}
\date{}

\begin{document}
\maketitle

\section*{Verdict}

\[
\boxed{\mathrm{final\_verdict}=V_{\mathrm{dichotomy\_formalized}}.}
\]

This step formalizes the Framework RH-Carrier Dichotomy as a candidate
foundational typed condition.  It combines the CRE foreclosure taxonomy
(\(\CRCFT\)) with the native-closure behavior observed in proved non-CRE RH
analogs.

\begin{theorem}[Framework RH-Carrier Dichotomy]
Let \(C\) be a typed primary carrier supporting a residual \(\XiC\) aimed at
closing an RH-analogous statement for an \(L\)-function or zeta function
\(L_C\).

\begin{enumerate}
\item If \(C\) is classically-RH-equivalent, then \(C\) exhibits \(\CRCFT\) in
at least one of the modes
\[
\CRCFT\text{-}\CTMT,\qquad
\CRCFT\text{-TE},\qquad
\CRCFT\text{-BF}.
\]
The framework output is the typed external-content interface naming the
precise obstruction.
\item If \(C\) is non-\(\CRE\), and the RH analog for \(L_C\) has been proved
with a complete native spectral, cohomological, or algebraic ledger, then the
framework yields native closure
\[
\XiC=0
\]
through that ledger.
\end{enumerate}
\end{theorem}

\section*{Instance Table}

\begin{center}
\begin{tabular}{p{0.23\linewidth}p{0.12\linewidth}p{0.19\linewidth}p{0.34\linewidth}p{0.08\linewidth}}
\toprule
Carrier & CRE & Outcome & Mode or mechanism & Step\\
\midrule
Burnol/Sonine pulled-evaluator & yes & \(\CRCFT\) foreclosure
& \(\CTMT\): A/B stuck at \(\kappa\); C foreclosed by nonzero \(L_{\rho,k}\)
& 145,162,173--180\\
Hecke \(L\)-function carrier & yes & \(\CRCFT\) foreclosure
& BF: H6 V-NC ledger non-comparability
& 167--168\\
de Branges \(H(E_{\rm RH})\) & yes & \(\CRCFT\) foreclosure
& TE: closure equivalent to RH; Conrey--Li survival
& 181\\
Hilbert--Polya / Berry--Keating & yes & \(\CRCFT\) foreclosure
& BF for standard \(xp\); TE for modified BK
& 182\\
Connes adelic / NCG & yes & \(\CRCFT\) foreclosure
& TE: full trace-formula closure equivalent to RH
& 184\\
Selberg/Maass & no & native closure
& \(\Xi_\Gamma=0\) by full Selberg trace ledger
& 185\\
Weil/Deligne function field & no & native closure
& \(\Xi_{\rm WD}=0\) by Deligne purity and Grothendieck--Lefschetz
& 186\\
\bottomrule
\end{tabular}
\end{center}

\section*{Consequence Theorem}

Applied to RH-analogous carriers, the framework discipline produces:
\begin{itemize}
\item typed external-content interfaces for CRE carriers, naming the exact
\(\CRCFT\) obstruction;
\item native closure certificates for non-CRE carriers whose RH analog is
proved with a complete native ledger.
\end{itemize}

In particular, attempting Riemann RH through any tested CRE carrier produces
only a \(\CRCFT\)-typed foreclosure under inherited records.  A live route to
Riemann RH would require either a non-CRE carrier whose native closure
transports to Riemann RH, which is currently absent, or an external theorem
resolving a \(\CRCFT\)-stuck record such as the Burnol \(\kappa\) theorem.

\begin{conjecture}[Framework RH-Carrier Dichotomy Coverage]
Every RH-analogous typed carrier is classified by the dichotomy into exactly
one of \(\CRCFT\) foreclosure or native closure.  No third class exists.
\end{conjecture}

This is a conjecture, not a theorem.  It is supported by seven carriers:
five CRE foreclosure instances and two non-CRE native closures.  It is
refutable by exhibiting an RH-analogous carrier outside both classes.

\section*{Relationship to CTMT and CRCFT}

\(\CTMT\), formalized at Step 180, is the post-reduction matrix-element
terminus condition.  \(\CRCFT\), formalized at Step 183, contains \(\CTMT\) as
one foreclosure mode and adds target-equivalence and bridge-failure.  The
dichotomy is the full carrier-level picture: \(\CRCFT\) is the CRE branch,
while proved non-CRE analogs close natively.

\section*{Corpus Candidate}

The dichotomy is flagged for future foundational-corpus inclusion in the
standalone note \texttt{dichotomy\_typed\_condition.md}.  Corpus integration
is not performed in this step.

\section*{Nonclaim Boundary}

The dichotomy does not prove RH, does not close any Riemann-zeta residual, and
does not transfer Selberg or Weil/Deligne closure to Riemann RH.  The coverage
statement is conjectural.

\end{document}
